\section{Related Literature}\label{sec:literature}

The paper combines nonlinear service pricing, commitment, sunk integration, and dynamic tariff choice. These ingredients are established separately in prior work; the contribution is the feedback from an architecture-induced installed state back into the profitability of that same architecture.

\citet{sundararajan2004} studies unlimited-usage and usage-based pricing when usage pricing can be costly. \citet{ladas2022} compare selling with pay-per-use service provision. We therefore make no novelty claim for fixed-versus-usage pricing or metering costs.

Dynamic commitment and adoption problems are also well established. \citet{penmetsa2015} study subscription pricing with strategic customers. \citet{gans2012} studies application pricing when platform access precedes later pricing. \citet{mutherswismer2022} show that tariff form can mitigate hold-up created by sunk seller participation. \citet{minryu2025} analyze investment that is sunk before a monopolistic supplier offers two-part tariffs. Hence future pricing effects on sunk adoption and tariff-form commitment are not new here.

Mixed pricing choices are likewise not a standalone contribution. \citet{wanghu2014} obtain a unique mixed-strategy equilibrium in a model of committed versus contingent pricing under competition. In the present paper, mixing is instead the equilibrium resolution of the two-threshold self-consistency gap.

Recent AI and cloud work further rules out application-based novelty. \citet{bergemannbonattismolin2025} study LLM pricing with token costs, fine-tuning, heterogeneous users, and two-part-tariff implementation. \citet{bergemannwang2025} study dynamic cloud contracts and committed-spend implementation. \citet{bhaskaran2026} analyze costly digital services including AI/cloud applications, while \citet{bichuchyaish2026} study costly current use and future quality in generative AI.

Our claim is therefore narrow. On strict $\mathcal R^+$, expected metering lowers high-use integration, while a larger high-use installed state raises the provider's metering gain. This reciprocal state feedback splits one fixed-state architecture threshold into two self-consistency thresholds. Corollary~\ref{cor:collapse} removes the feedback and collapses the split. Generic hold-up, nonlinear pricing, mixed strategies, AI as an application label, and model improvement are not claimed as standalone contributions.
